\documentclass[14pt]{article}
\usepackage[14pt]{extsizes}
\usepackage[utf8]{inputenc}
\usepackage[T2A]{fontenc}
\usepackage[english, russian]{babel}
\usepackage[a4paper,left=30mm, right=10mm, top=20mm, bottom=20mm]{geometry}
\usepackage{indentfirst, setspace}
\setlength{\parindent}{1.25cm}
\usepackage{tabularx, multirow}
\usepackage[normalem]{ulem}
\usepackage[style=russian]{csquotes}
\usepackage{graphicx,wrapfig}
\fontsize{12}{12}\selectfont
\usepackage{amsmath,amsfonts,amssymb,amsthm} % Добавлен amsthm
\usepackage{mathtools}
\usepackage{listings}
\usepackage{caption}
\usepackage{enumitem} % ДОБАВЛЕН ПАКЕТ ДЛЯ [label=...]
\renewcommand{\labelenumii}{\arabic{enumi}.\arabic{enumii}.}
\renewcommand{\labelenumiii}{\arabic{enumi}.\arabic{enumii}.\arabic{enumiii}.}
\renewcommand{\labelenumiv}{\arabic{enumi}.\arabic{enumii}.\arabic{enumiii}.\arabic{enumiv}.}
\DeclareCaptionLabelSeparator{custom}{ --- }
\captionsetup{labelsep=custom}
\usepackage{pgfplots}
\usepackage{pdfpages}
\pgfplotsset{compat=1.9}
\usepackage{xcolor}
\usepackage{hyperref}
\definecolor{linkcolor}{HTML}{000000}
\definecolor{urlcolor}{HTML}{000000}
\hypersetup{pdfstartview=FitH, linkcolor=linkcolor, urlcolor=urlcolor, colorlinks=true}
\captionsetup[figure]{name=Рисунок}

% ---------- Центрирование заголовков задач ----------
\usepackage{titlesec}
\titleformat{\subsection}[block]
{\centering\normalfont\large\bfseries}{}{0pt}{}
% ----------------------------------------------------

\newtheorem*{solution}{Решение}

\begin{document}
	
	% ============ ВАРИАНТ 1 ============
	\section{Вариант 1}
	
	\subsection*{Задача 1}
	Пусть $\mathbb{H}$ — тело кватернионов. Рассмотрим множество матриц
	\[
	H' = \left\{ 
	\begin{pmatrix} 
		u & v \\ 
		-\bar{v} & \bar{u} 
	\end{pmatrix} \;\bigg|\; u, v \in \mathbb{C} 
	\right\}.
	\]
	Это множество является подкольцом кольца $M_2(\mathbb{C})$.
	
	Определим отображение $f: H' \to \mathbb{H}$ формулой:
	\[
	f \begin{pmatrix} 
		a+bi & c+di \\ 
		-c+di & a-bi 
	\end{pmatrix} 
	= a + bi + cj + dk.
	\]
	
	\begin{enumerate}
		\item \textbf{Биективность:} Любой кватернион $q = a + bi + cj + dk \in \mathbb{H}$ однозначно определяется набором вещественных коэффициентов $(a, b, c, d)$. Матрица из подкольца $H'$ при подстановке комплексных чисел $u = a + bi$ и $v = c + di$ также однозначно задается этой же четверкой чисел. Соответствие взаимно однозначно.
		
		\item \textbf{Аддитивность:} Так как сложение матриц и сложение кватернионов определены покомпонентно и линейны по переменным $a, b, c, d$, то равенство $f(A + B) = f(A) + f(B)$ выполняется автоматически.
		
		\item \textbf{Мультипликативность:} Разложим матрицу из $H'$ по вещественному базису $A = a \cdot E + b \cdot I + c \cdot J + d \cdot K$, где $E, I, J, K$ — базисные матрицы. Отображение $f$ переводит их в стандартные единицы кватернионов: $f(E)=1$, $f(I)=i$, $f(J)=j$, $f(K)=k$. Прямым перемножением матриц проверяются соотношения:
		\[
		I^2 = J^2 = K^2 = -E, \quad I \cdot J = K = -J \cdot I, \quad J \cdot K = I = -K \cdot J, \quad K \cdot I = J = -I \cdot K.
		\]
		Эта таблица умножения полностью тождественна алгебре кватернионов, откуда $f(A \cdot B) = f(A) \cdot f(B)$.
	\end{enumerate}
	
	\textbf{Вывод:} Отображение $f$ — изоморфизм колец.
	
	\subsection*{Задача 2}
	Пусть $R = \{ a + b\sqrt{3} \mid a, b \in \mathbb{Z} \}$, $J = \{ a + b\sqrt{3} \mid a, b \in 2\mathbb{Z} \}$. Доказать, что элементы $a + b\sqrt{3}$ и $c + d\sqrt{3}$ принадлежат одному и тому же классу вычетов по модулю $J$ тогда и только тогда, когда $a \equiv c \pmod{2}$ и $b \equiv d \pmod{2}$.
	
	\begin{solution}
		Элементы принадлежат одному классу вычетов по модулю идеала $J$ тогда и только тогда, когда их разность лежит в $J$.
		
		Вычислим разность:
		\[
		(a + b\sqrt{3}) - (c + d\sqrt{3}) = (a - c) + (b - d)\sqrt{3}.
		\]
		
		По определению идеала $J$, элемент $(a - c) + (b - d)\sqrt{3} \in J$ тогда и только тогда, когда коэффициенты при базисных элементах являются четными числами:
		\begin{enumerate}
			\item $(a - c) \in 2\mathbb{Z} \iff a - c \equiv 0 \pmod{2} \iff a \equiv c \pmod{2}$;
			\item $(b - d) \in 2\mathbb{Z} \iff b - d \equiv 0 \pmod{2} \iff b \equiv d \pmod{2}$.
		\end{enumerate}
		
		Утверждение полностью доказано.
	\end{solution}
	
	\subsection*{Задача 3}
	Пусть $f: \mathbb{C} \to \mathbb{R}$ задано формулой $f(a + bi) = a$. Доказать, что при любых $z_1, z_2 \in \mathbb{C}$ выполнено $f(z_1 + z_2) = f(z_1) + f(z_2)$, но существуют такие $z_1, z_2$, что $f(z_1 z_2) \neq f(z_1)f(z_2)$.
	
	\begin{solution}
		\begin{enumerate}
			\item \textbf{Проверка аддитивности:}
			Пусть $z_1 = a_1 + b_1 i$, $z_2 = a_2 + b_2 i$. Тогда $z_1 + z_2 = (a_1 + a_2) + (b_1 + b_2)i$.
			\[
			f(z_1 + z_2) = a_1 + a_2.
			\]
			С другой стороны, $f(z_1) + f(z_2) = a_1 + a_2$. Аддитивность доказана.
			
			\item \textbf{Опровержение мультипликативности (контрпример):}
			Пусть $z_1 = i$, $z_2 = i$. Тогда $z_1 \cdot z_2 = i^2 = -1$.
			Вычислим значения отображения:
			\[
			f(z_1 z_2) = f(-1) = -1.
			\]
			\[
			f(z_1) = f(i) = 0, \quad f(z_2) = f(i) = 0 \implies f(z_1) \cdot f(z_2) = 0 \cdot 0 = 0.
			\]
			Очевидно, что $-1 \neq 0$, следовательно $f(z_1 z_2) \neq f(z_1)f(z_2)$.
		\end{enumerate}
		Отображение сохраняет сложение, но не сохраняет умножение.
	\end{solution}
	
	\subsection*{Задача 3.5}
	Пусть отображение $f: M_2(\mathbb{C}) \to \mathbb{C}$ задано формулой $f(A) = \det A$. Доказать, что при любых $A, B \in M_2(\mathbb{C})$ выполнено $f(AB) = f(A)f(B)$, но существуют такие матрицы $A$ и $B$, что $f(A + B) \neq f(A) + f(B)$.
	
	\begin{solution}
		\begin{enumerate}
			\item \textbf{Мультипликативность:} По фундаментальному свойству определителя матрицы, для любых квадратных матриц $A$ и $B$ справедливо равенство $\det(AB) = \det(A) \cdot \det(B)$. Следовательно, $f(AB) = f(A)f(B)$ выполняется всегда.
			
			\item \textbf{Опровержение аддитивности (контрпример):}
			Рассмотрим единичные матрицы:
			\[
			A = \begin{pmatrix} 1 & 0 \\ 0 & 1 \end{pmatrix}, \quad B = \begin{pmatrix} 1 & 0 \\ 0 & 1 \end{pmatrix}.
			\]
			Тогда $\det(A) = 1$, $\det(B) = 1$, значит $f(A) + f(B) = 1 + 1 = 2$.
			
			Сложим матрицы:
			\[
			A + B = \begin{pmatrix} 2 & 0 \\ 0 & 2 \end{pmatrix}.
			\]
			Тогда $f(A + B) = \det(A + B) = 2 \cdot 2 - 0 \cdot 0 = 4$.
			
			Поскольку $4 \neq 2$, свойство аддитивности в общем случае не выполняется.
		\end{enumerate}
	\end{solution}
	
	\subsection*{Задача 4}
	Доказать, что число $\alpha = \sqrt{2 - \sqrt{3}} + 1$ является алгебраическим над $\mathbb{Q}$.
	
	\begin{solution}
		Число называется алгебраическим над $\mathbb{Q}$, если оно является корнем ненулевого многочлена с рациональными коэффициентами. Выразим $x = \alpha$ и избавимся от иррациональностей:
		\[
		x = \sqrt{2 - \sqrt{3}} + 1 \implies x - 1 = \sqrt{2 - \sqrt{3}}.
		\]
		Возведем обе части в квадрат:
		\[
		(x - 1)^2 = 2 - \sqrt{3} \implies x^2 - 2x + 1 = 2 - \sqrt{3} \implies x^2 - 2x - 1 = -\sqrt{3}.
		\]
		Снова возведем в квадрат, чтобы уничтожить $\sqrt{3}$:
		\[
		(x^2 - 2x - 1)^2 = (-\sqrt{3})^2 \implies (x^2 - 2x - 1)^2 = 3.
		\]
		Раскроем левую часть:
		\begin{align*}
			(x^2 - 2x - 1)(x^2 - 2x - 1) &= x^4 - 2x^3 - x^2 - 2x^3 + 4x^2 + 2x - x^2 + 2x + 1 \\
			&= x^4 - 4x^3 + 2x^2 + 4x + 1.
		\end{align*}
		Приравняем к 3:
		\[
		x^4 - 4x^3 + 2x^2 + 4x + 1 = 3 \implies x^4 - 4x^3 + 2x^2 + 4x - 2 = 0.
		\]
		Многочлен $f(x) = x^4 - 4x^3 + 2x^2 + 4x - 2$ имеет рациональные коэффициенты, и $f(\alpha) = 0$. Число $\alpha$ является алгебраическим.
	\end{solution}
	
	\subsection*{Задача 5}
	Представьте многочлен $f = x_1^4 x_2 + x_1^4 x_3 + x_1 x_2^4 + x_1 x_3^4 + x_2^4 x_3 + x_2 x_3^4$ в виде многочлена от элементарных симметрических многочленов $\sigma_1, \sigma_2, \sigma_3$.
	
	\begin{solution}
		Элементарные симметрические многочлены для 3 переменных:
		\begin{align*}
			\sigma_1 &= x_1 + x_2 + x_3, \\
			\sigma_2 &= x_1 x_2 + x_1 x_3 + x_2 x_3, \\
			\sigma_3 &= x_1 x_2 x_3.
		\end{align*}
		
		Исходный многочлен можно записать через орбиту монома: $f = \sum_{\text{sym}} x_1^4 x_2$. 
		Старший член $f$ в лексикографическом порядке равен $x_1^4 x_2$. Ему соответствует комбинация $\sigma_1^3 \sigma_2$: старший член $\sigma_1^3 \sigma_2$ равен $(x_1)^3 \cdot (x_1 x_2) = x_1^4 x_2$.
		
		Вычтем её и продолжим алгоритм (метод Гаусса для симметрических многочленов):
		\[
		f = \sigma_1^2 \sigma_2^2 - 2\sigma_2^3 - 3\sigma_1^3 \sigma_3 + 5\sigma_1 \sigma_2 \sigma_3.
		\]
	\end{solution}
	
	\newpage
	
	% ============ ВАРИАНТ 2 ============
	\section{Вариант 2}
	
	\subsection*{Задача 1}
	Докажите, что центром тела кватернионов является поле действительных чисел $\mathbb{R}$.
	
	\begin{solution}
		Центр $Z(\mathbb{H})$ состоит из элементов $q = a + bi + cj + dk$, коммутирующих со всеми элементами тела. В частности, $q$ должен коммутировать с мнимыми единицами $i$ и $j$.
		
		\begin{enumerate}
			\item \textbf{Коммутирование с $i$:}
			\begin{align*}
				q \cdot i &= (a + bi + cj + dk)i = ai - b - ck + dj, \\
				i \cdot q &= i(a + bi + cj + dk) = ai - b + ck - dj.
			\end{align*}
			$q \cdot i = i \cdot q \implies -ck + dj = ck - dj \implies 2ck = 0, \; 2dj = 0 \implies c = 0, \; d = 0$. Элемент сужается до вида $q = a + bi$.
			
			\item \textbf{Коммутирование с $j$:}
			\begin{align*}
				q \cdot j &= (a + bi)j = aj + bk, \\
				j \cdot q &= j(a + bi) = aj - bk.
			\end{align*}
			$q \cdot j = j \cdot q \implies bk = -bk \implies 2bk = 0 \implies b = 0$.
		\end{enumerate}
		
		Остается $q = a$, где $a \in \mathbb{R}$. Действительные числа коммутируют со всеми кватернионами по определению.
	\end{solution}
	
	\subsection*{Задача 2}
	Согласно задаче 2 (Вариант 1), класс вычетов элемента $a + b\sqrt{3}$ однозначно определяется остатками чисел $a$ и $b$ от деления на $2$.
	
	\begin{solution}
		В кольце целых чисел $\mathbb{Z}$ по модулю $2$ существуют два возможных остатка: $\{0, 1\}$. Следовательно, для пары коэффициентов $(a, b)$ существует ровно $2 \times 2 = 4$ различные комбинации остатков:
		\begin{enumerate}
			\item $a \equiv 0$, $b \equiv 0 \implies$ класс $0 + 0\sqrt{3} = u_0 = \bar{0}$;
			\item $a \equiv 1$, $b \equiv 0 \implies$ класс $1 + 0\sqrt{3} = u_1 = \bar{1}$;
			\item $a \equiv 0$, $b \equiv 1 \implies$ класс $0 + 1\sqrt{3} = u_2 = \overline{\sqrt{3}}$;
			\item $a \equiv 1$, $b \equiv 1 \implies$ класс $1 + 1\sqrt{3} = u_3 = \overline{1+\sqrt{3}}$.
		\end{enumerate}
		Все эти 4 класса попарно различны, так как их представители имеют разные остатки коэффициентов по модулю $2$, и в сумме они исчерпывают всё кольцо $R$. Фактор-кольцо $R/J$ состоит из 4 указанных элементов.
	\end{solution}
	
	\subsection*{Задача 3}
	Пусть $R$ — область целостности характеристики $p$. Доказать, что:
	\begin{enumerate}
		\item при $a \in R$, $a \neq 0$ и $k, l \in \mathbb{Z}$: $ka = la \iff k \equiv l \pmod{p}$;
		\item при $a, b \in R$ и $k \not\equiv 0 \pmod{p}$: $ka = kb \iff a = b$.
	\end{enumerate}
	
	\begin{solution}
		Характеристика области целостности равна $p$ (всегда простое число). Это означает, что $p \cdot 1 = 0$, и $p$ — наименьшее натуральное число с таким свойством.
		
		\textbf{Пункт 1:} $ka = la \iff ka - la = 0 \iff (k - l)a = 0$.
		
		Поскольку $R$ — область целостности и $a \neq 0$, равенство $(k - l)a = 0$ равносильно тому, что сам скаляр $(k - l)$ в кольце равен нулю, то есть кратен характеристике кольца $p$. Значит, $(k - l) \vdots p$, что эквивалентно $k \equiv l \pmod{p}$.
		
		\textbf{Пункт 2:} $ka = kb \iff ka - kb = 0 \iff k(a - b) = 0$.
		
		Так как $k \not\equiv 0 \pmod{p}$, элемент $k \cdot 1 \neq 0$ в кольце $R$. Поскольку в области целостности нет делителей нуля, из произведения $k \cdot (a - b) = 0$ при $k \neq 0$ следует, что второй множитель обязан быть нулем: $a - b = 0 \iff a = b$.
	\end{solution}
	
	\subsection*{Задача 4}
	Докажите изоморфизм полей $\mathbb{Q}(\sqrt[3]{5})$ и $\mathbb{Q}[x]/(x^3 - 5)$.
	
	\begin{solution}
		Рассмотрим отображение подстановки (эвалюации) $\psi: \mathbb{Q}[x] \to \mathbb{Q}(\sqrt[3]{5})$, заданное формулой $\psi(f(x)) = f(\sqrt[3]{5})$.
		
		\begin{enumerate}
			\item \textbf{Гомоморфизм и сюръективность:} Отображение $\psi$ является гомоморфизмом колец. Так как поле $\mathbb{Q}(\sqrt[3]{5})$ по определению состоит из элементов вида $a + b\sqrt[3]{5} + c(\sqrt[3]{5})^2$, где $a,b,c \in \mathbb{Q}$, то любой элемент поля является образом многочлена $a + bx + cx^2 \in \mathbb{Q}[x]$. Следовательно, отображение сюръективно ($\operatorname{Im}(\psi) = \mathbb{Q}(\sqrt[3]{5})$).
			
			\item \textbf{Нахождение ядра:} Многочлен $p(x) = x^3 - 5$ неприводим над $\mathbb{Q}$ по критерию Эйзенштейна (с простым $p=5$). Число $\sqrt[3]{5}$ является его корнем. Следовательно, $p(x)$ — минимальный многочлен для $\sqrt[3]{5}$ над $\mathbb{Q}$. Ядро гомоморфизма $\operatorname{Ker}(\psi)$ состоит из всех многочленов, зануляющихся в точке $\sqrt[3]{5}$, а значит, делящихся на минимальный многочлен. То есть $\operatorname{Ker}(\psi) = \langle x^3 - 5 \rangle$.
		\end{enumerate}
		По основной теореме о гомоморфизме колец, фактор-кольцо по ядру изоморфно образу:
		\[
		\mathbb{Q}[x]/(x^3 - 5) \cong \mathbb{Q}(\sqrt[3]{5}).
		\]
	\end{solution}
	
	\subsection*{Задача 5}
	При каких значениях параметра $\lambda$ многочлены $f(x) = x^3 - \lambda x + 2$ и $g(x) = x^2 + \lambda x + 2$ имеют общий корень?
	
	\begin{solution}
		Пусть $x_0$ — общий корень многочленов. Тогда $f(x_0) = 0$ и $g(x_0) = 0$:
		\begin{align}
			x_0^3 - \lambda x_0 + 2 &= 0, \\
			x_0^2 + \lambda x_0 + 2 &= 0.
		\end{align}
		
		Сложим эти два уравнения, чтобы избавиться от параметра $\lambda$:
		\[
		x_0^3 + x_0^2 + 4 = 0.
		\]
		
		Найдем целые корни этого уравнения среди делителей свободного члена $4$ ($\pm 1, \pm 2, \pm 4$).
		Проверим $x_0 = -2$: $(-2)^3 + (-2)^2 + 4 = -8 + 4 + 4 = 0$. Корень найден!
		
		Разделив многочлен на $(x_0 + 2)$ уголком или по схеме Горнера, получим разложение: 
		\[
		(x_0 + 2)(x_0^2 - x_0 + 2) = 0.
		\]
		
		Случай 1 (вещественный корень): Пусть $x_0 = -2$. Подставим его во второе уравнение (2) для поиска параметра $\lambda$:
		\[
		(-2)^2 + \lambda(-2) + 2 = 0 \implies 4 - 2\lambda + 2 = 0 \implies 6 = 2\lambda \implies \lambda = 3.
		\]
		
		Случай 2 (комплексные корни): Корни из второй скобки удовлетворяют уравнению $x_0^2 - x_0 + 2 = 0$, откуда $x_0^2 + 2 = x_0$. Подставим это во второе уравнение системы (2):
		\[
		(x_0^2 + 2) + \lambda x_0 = 0 \implies x_0 + \lambda x_0 = 0 \implies x_0(1 + \lambda) = 0.
		\]
		Поскольку комплексные корни $x_0 \neq 0$, сокращаем на $x_0$ и получаем $\lambda = -1$.
		
		\textbf{Вывод:} Многочлены имеют общие корни при $\lambda = 3$ (корень $x_0 = -2$) и при $\lambda = -1$.
	\end{solution}
	
	\newpage
	
	% ============ ВАРИАНТ 3 ============
	\section{Вариант 3}
	
	\subsection*{Задача 1}
	Пусть $L$ — подтело тела $D$. Дать определение левого линейного пространства над $L$ и доказать, что $D$ таковым является.
	
	\begin{solution}
		\textbf{Определение.} Левым линейным пространством над телом (или полем) $L$ называется абелева группа $(V,+)$ с заданной операцией внешнего левого умножения на скаляр $L \times V \to V$, удовлетворяющей для любых $\alpha,\beta \in L$ и $u,v \in V$ следующим четырем аксиомам:
		\begin{enumerate}[label=(\arabic*)]
			\item $\alpha(u+v) = \alpha u + \alpha v$ (дистрибутивность скаляра относительно суммы векторов);
			\item $(\alpha+\beta)v = \alpha v + \beta v$ (дистрибутивность вектора относительно суммы скаляров);
			\item $(\alpha\beta)v = \alpha(\beta v)$ (ассоциативность умножения на скаляры);
			\item $1_L \cdot v = v$ (унитарность, где $1_L$ — единица тела $L$).
		\end{enumerate}
		
		\textbf{Доказательство для $V=D$:} 
		Пусть в качестве множества векторов $V$ выступает само тело $D$ со своей внутренней операцией сложения. В качестве скаляров берутся элементы подтела $L \subseteq D$. Внешнее умножение скаляра на вектор зададим как обычное умножение элементов слева внутри самого кольца $D$.
		
		Тогда все аксиомы линейного пространства автоматически следуют из базовых аксиом кольца $D$:
		\begin{enumerate}[label=(\arabic*)]
			\item Первая аксиома выполняется в силу свойства левой дистрибутивности умножения относительно сложения в теле $D$.
			\item Вторая аксиома выполняется в силу свойства правой дистрибутивности в теле $D$.
			\item Третья аксиома строго выполняется благодаря ассоциативности умножения внутри $D$.
			\item Так как $L$ является подтелом, его единица $1_L$ совпадает с единицей $1_D$ кольца $D$. По определению единичного элемента в кольце $1_D \cdot v = v$ для любого $v \in D$.
		\end{enumerate}
		Следовательно, тело $D$ образует корректное левое линейное пространство над своим подтелом $L$.
	\end{solution}
	
	\subsection*{Задача 2}
	Пусть $K$ — числовое поле, $R=K[x]$ — кольцо многочленов, $J=(x-1)R$ — идеал. Доказать, что $f(x),g(x)$ принадлежат одному классу по модулю $J$ тогда и только тогда, когда $f(1)=g(1)$.
	
	\begin{solution}
		По определению фактор-кольца, многочлены $f(x)$ и $g(x)$ лежат в одном и том же смежном классе вычетов по модулю идеала $J = \langle x-1 \rangle$ тогда и только тогда, когда их разность полностью принадлежит этому идеалу:
		\[
		f(x) \equiv g(x) \pmod J \iff f(x) - g(x) \in J.
		\]
		Принадлежность многочлена главному идеалу $J = (x-1)R$ означает, что этот многочлен делится нацело на образующий элемент $(x-1)$. Иными словами, существует многочлен $q(x) \in K[x]$ такой, что:
		\[
		f(x) - g(x) = (x-1)q(x).
		\]
		\textbf{Необходимость ($\Rightarrow$):} Если элементы лежат в одном классе, то выполнено $f(x) - g(x) = (x-1)q(x)$. Подставим в это тождество значение $x = 1$:
		\[
		f(1) - g(1) = (1-1)q(1) = 0 \cdot q(1) = 0 \implies f(1) = g(1).
		\]
		\textbf{Достаточность ($\Leftarrow$):} Пусть известно, что $f(1) = g(1)$. Рассмотрим новый многочлен $h(x) = f(x) - g(x)$. Так как $h(1) = f(1) - g(1) = 0$, то по теореме Безу число $x=1$ является корнем многочлена $h(x)$. Следовательно, $h(x)$ делится на $(x-1)$ без остатка, то есть $h(x) \in \langle x-1 \rangle = J$. А значит, $f$ и $g$ лежат в одном классе.
	\end{solution}
	
	\subsection*{Задача 3}
	Пусть $R$ — область целостности характеристики $p$. Доказать, что для любых $a,b\in R$
	\[
	(a+b)^p=a^p+b^p,\qquad (a-b)^p=a^p-b^p.
	\]
	
	\begin{solution}
		Характеристика области целостности $R$ равна $p$, что означает $p \cdot 1_R = 0$. При этом $p$ обязано быть простым числом (иначе в $R$ существовали бы делители нуля).
		
		Рассмотрим разложение степени суммы по формуле бинома Ньютона (которая справедлива в любом коммутативном кольце):
		\[
		(a+b)^p = a^p + \sum_{k=1}^{p-1}\binom{p}{k}a^{p-k}b^k + b^p.
		\]
		Проанализируем промежуточные биномиальные коэффициенты $\binom{p}{k} = \frac{p!}{k!(p-k)!}$ для $0 < k < p$. В числителе присутствует множитель $p$. Знаменатель является произведением натуральных чисел, строго меньших простого числа $p$. Следовательно, знаменатель не может содержать делитель $p$ и не может его сократить. Значит, число $\binom{p}{k}$ нацело делится на $p$.
		
		Так как любое кратное $p$ в кольце $R$ равно нулю ($p \cdot c = 0$), все промежуточные слагаемые суммы обращаются в нуль:
		\[
		(a+b)^p = a^p + 0 + b^p = a^p + b^p.
		\]
		
		Для разности применим только что доказанную формулу к элементам $a$ и $(-b)$:
		\[
		(a-b)^p = (a+(-b))^p = a^p + (-b)^p.
		\]
		Поскольку $p$ — простое число, возможны два случая:
		\begin{enumerate}
			\item Если $p=2$, то характеристика кольца равна 2, а значит $1 = -1$ (ведь $1+1=0$). Тогда $(-b)^2 = b^2$. Но и $a^2 + b^2 = a^2 - b^2$, так как плюсы и минусы в таком кольце эквивалентны.
			\item Если $p > 2$, то $p$ — нечетное простое число. Нечетная степень отрицательного элемента дает минус: $(-b)^p = -b^p$. Подставляя это, получаем строго $(a-b)^p = a^p - b^p$.
		\end{enumerate}
		В обоих случаях тождество выполняется.
	\end{solution}
	
	\subsection*{Задача 4}
	Пусть $a,b,r\in \mathbb Q$ и $\sqrt r \notin \mathbb Q$. Доказать, что $\alpha = a+b\sqrt r$ — алгебраическое число, и найти его минимальный многочлен над $\mathbb Q$.
	
	\begin{solution}
		Обозначим $x = \alpha$. Произведем алгебраические преобразования, чтобы построить аннулирующий многочлен с коэффициентами из $\mathbb Q$:
		\[
		x = a + b\sqrt{r} \implies x - a = b\sqrt{r}.
		\]
		Возведем обе части уравнения в квадрат:
		\[
		(x - a)^2 = (b\sqrt{r})^2 \implies x^2 - 2ax + a^2 = b^2r \implies x^2 - 2ax + (a^2 - b^2r) = 0.
		\]
		Получен многочлен $f(x) = x^2 - 2ax + (a^2 - b^2r)$. Так как числа $a, b, r$ рациональные, коэффициенты $f(x)$ также принадлежат $\mathbb Q$. И по построению $f(\alpha) = 0$. Следовательно, $\alpha$ — алгебраическое число над $\mathbb Q$.
		
		Теперь докажем, что $f(x)$ является минимальным многочленом $m_\alpha(x)$. Минимальный многочлен должен быть унитарным (старший коэффициент 1) и неприводимым над полем $\mathbb Q$. Наш многочлен $f(x)$ является унитарным и имеет степень 2.
		Если бы $f(x)$ был приводим над $\mathbb Q$, то он бы раскладывался на два линейных множителя с рациональными коэффициентами, а значит, имел бы исключительно рациональные корни. 
		Однако его корнями являются числа $x_{1,2} = a \pm b\sqrt{r}$. По условию задачи $b \neq 0$ (иначе бы корень не фигурировал в условии $\sqrt{r} \notin \mathbb{Q}$ как значимый) и $\sqrt{r} \notin \mathbb{Q}$. Сумма рационального и иррационального чисел иррациональна, поэтому корни не лежат в $\mathbb Q$. Следовательно, многочлен неприводим.
		
		\textbf{Вывод:} $m_\alpha(x) = x^2 - 2ax + (a^2 - b^2r)$.
	\end{solution}
	
	\subsection*{Задача 5}
	Пользуясь свойствами результанта, решите систему
	\[
	\begin{cases}
		4x^2 - 7xy + 13x + y^2 - 2y - 3 = 0,\\
		9x^2 - 14xy + 28x + y^2 - 4y - 5 = 0.
	\end{cases}
	\]
	
	\begin{solution}
		Рассмотрим уравнения системы как два многочлена $f(y)$ и $g(y)$ относительно переменной $y$, коэффициентами которых являются выражения от $x$:
		\begin{align*}
			f(y) &= y^2 - (7x + 2)y + (4x^2 + 13x - 3), \\
			g(y) &= y^2 - (14x + 4)y + (9x^2 + 28x - 5).
		\end{align*}
		По теореме о результанте, система имеет решение по переменной $y$ тогда и только тогда, когда результант этих многочленов равен нулю: $\operatorname{Res}_y(f, g) = 0$.
		
		Воспользуемся важнейшим свойством результанта $\operatorname{Res}(A, B) = \operatorname{Res}(A, B - A)$. Пусть $h(y) = g(y) - f(y)$:
		\[
		h(y) = -(7x + 2)y + (5x^2 + 15x - 2).
		\]
		Теперь необходимо вычислить результант квадратного многочлена $f(y) = y^2 + py + q$ и линейного многочлена $h(y) = Ay + B$, где:
		\[
		p = -(7x+2), \quad q = 4x^2+13x-3, \quad A = -(7x+2), \quad B = 5x^2+15x-2.
		\]
		По формуле результанта для степеней 2 и 1:
		\[
		\operatorname{Res}(f, h) = A^2 f\left(-\frac{B}{A}\right) = B^2 - p A B + q A^2.
		\]
		Подставляем значения:
		\[
		\operatorname{Res}_y = B^2 - (-(7x+2))(-(7x+2))B + q(7x+2)^2 = B^2 - (7x+2)^2 B + (7x+2)^2 q.
		\]
		Выносим общий множитель:
		\[
		\operatorname{Res}_y = B^2 - (7x+2)^2 (B - q).
		\]
		Вычислим разность $(B - q)$:
		\[
		B - q = (5x^2 + 15x - 2) - (4x^2 + 13x - 3) = x^2 + 2x + 1 = (x+1)^2.
		\]
		Тогда уравнение результанта принимает вид разности квадратов:
		\[
		(5x^2 + 15x - 2)^2 - ((7x+2)(x+1))^2 = 0.
		\]
		Раскладываем по формуле $U^2 - V^2 = (U-V)(U+V)$, предварительно раскрыв $(7x+2)(x+1) = 7x^2 + 9x + 2$:
		\begin{align*}
			U - V &= (5x^2 + 15x - 2) - (7x^2 + 9x + 2) = -2x^2 + 6x - 4 = -2(x^2 - 3x + 2) = -2(x-1)(x-2). \\
			U + V &= (5x^2 + 15x - 2) + (7x^2 + 9x + 2) = 12x^2 + 24x = 12x(x+2).
		\end{align*}
		Перемножая эти две скобки, получаем итоговый результант:
		\[
		\operatorname{Res}_y(f,g) = -24x(x-1)(x-2)(x+2) = 0.
		\]
		Корни данного уравнения: $x_1 = 0$, $x_2 = 1$, $x_3 = 2$, $x_4 = -2$.
		
		Для нахождения $y$ подставим найденные $x$ в линейное уравнение $h(y) = 0$, откуда $y = -B/A = \frac{5x^2 + 15x - 2}{7x + 2}$:
		\begin{itemize}
			\item При $x = 0$: $y = \frac{-2}{2} = -1 \implies (0, -1)$.
			\item При $x = 1$: $y = \frac{5 + 15 - 2}{7 + 2} = \frac{18}{9} = 2 \implies (1, 2)$.
			\item При $x = 2$: $y = \frac{20 + 30 - 2}{14 + 2} = \frac{48}{16} = 3 \implies (2, 3)$.
			\item При $x = -2$: $y = \frac{20 - 30 - 2}{-14 + 2} = \frac{-12}{-12} = 1 \implies (-2, 1)$.
		\end{itemize}
		Все четыре пары являются точными решениями исходной системы.
	\end{solution}
	
	\newpage
	
	% ============ ВАРИАНТ 4 ============
	\section{Вариант 4}
	
	\subsection*{Задача 1}
	Докажите, что $1,i,j,k$ — левый $\mathbb R$-базис тела кватернионов $\mathbb H$.
	
	\begin{solution}
		В алгебре тело кватернионов $\mathbb H$ по определению конструируется как четырехмерное векторное пространство над полем действительных чисел $\mathbb R$ с фиксированным базисом $\{1, i, j, k\}$ и определенной таблицей некоммутативного умножения для базисных элементов.
		
		Докажем, что эта система образует левый базис явно.
		\begin{enumerate}
			\item \textbf{Система образующих:} Любой кватернион $q\in\mathbb H$ по своему определению записывается в виде:
			\[
			q = a\cdot 1 + b\cdot i + c\cdot j + d\cdot k,\qquad a,b,c,d\in\mathbb R.
			\]
			Это означает, что любой элемент тела можно представить в виде левой линейной комбинации скаляров из $\mathbb R$ и элементов множества $\{1,i,j,k\}$. Следовательно, система порождает $\mathbb H$.
			
			\item \textbf{Линейная независимость:} Составим линейную комбинацию, приравненную к нулевому кватерниону:
			\[
			a\cdot 1 + b\cdot i + c\cdot j + d\cdot k = 0,\qquad a,b,c,d\in\mathbb R.
			\]
			По определению равенства кватернионов, два кватерниона равны тогда и только тогда, когда равны все их соответствующие компоненты. Поскольку справа стоит нулевой кватернион $0 + 0i + 0j + 0k$, то:
			\[
			a=0,\; b=0,\; c=0,\; d=0.
			\]
			Это означает, что линейная комбинация тривиальна, и базисные векторы линейно независимы над $\mathbb R$.
		\end{enumerate}
		Таким образом, $\{1,i,j,k\}$ строго образует левый $\mathbb R$-базис тела $\mathbb H$.
	\end{solution}
	
	\subsection*{Задача 2}
	Пусть $K$ — числовое поле, $R=K[x]$, $J=(x-1)R$. Доказать, что отображение
	\[
	R/J \to K,\qquad f(x)+J \mapsto f(1)
	\]
	является изоморфизмом.
	
	\begin{solution}
		Обозначим заданное отображение как $\Phi: R/J \to K$, $\Phi(f(x)+J) = f(1)$.
		
		\begin{enumerate}
			\item \textbf{Корректность отображения:} Покажем, что значение функции не зависит от выбора представителя смежного класса. Пусть $f(x)+J = g(x)+J$. Тогда $f(x) - g(x) \in J = \langle x-1 \rangle$. Следовательно, разность делится на $x-1$, то есть $f(x) - g(x) = (x-1)q(x)$. Подставляя $x=1$, получаем $f(1) - g(1) = 0 \implies f(1) = g(1)$. Отображение задано корректно.
			
			\item \textbf{Гомоморфизм колец:} Отображение сохраняет алгебраические операции сложения и умножения:
			\begin{align*}
				\Phi((f+J) + (g+J)) &= \Phi((f+g)+J) = (f+g)(1) = f(1)+g(1) = \Phi(f+J) + \Phi(g+J), \\
				\Phi((f+J) \cdot (g+J)) &= \Phi((fg)+J) = (fg)(1) = f(1)g(1) = \Phi(f+J) \cdot \Phi(g+J).
			\end{align*}
			
			\item \textbf{Сюръективность:} Для любого числа $c \in K$ существует многочлен-константа $h(x) = c$. Класс смежности этого многочлена отображается ровно в $c$: $\Phi(c+J) = c$. Образ отображения покрывает все поле $K$.
			
			\item \textbf{Инъективность:} Пусть $\Phi(f(x)+J) = 0$. По правилу отображения это значит, что $f(1) = 0$. По теореме Безу, если $1$ является корнем многочлена, то многочлен делится нацело на $(x-1)$. Это означает, что $f(x) \in (x-1)R = J$. Следовательно, класс $f(x)+J$ является нулевым классом вычетов $0+J$. Так как ядро состоит только из нуля, отображение инъективно.
		\end{enumerate}
		Отображение биективно и сохраняет операции, следовательно, это изоморфизм.
	\end{solution}
	
	\subsection*{Задача 3}
	Докажите, что в области целостности характеристики $p$ верно равенство
	\[
	(a+b)^p = a^p + b^p.
	\]
	
	\begin{solution}
		Характеристика $p$ кольца — это минимальное простое число, при котором любое кратное $p \cdot x = 0$.
		
		По классической формуле бинома Ньютона, которая верна в любом коммутативном кольце:
		\[
		(a+b)^p = a^p + \sum_{k=1}^{p-1}\binom{p}{k}a^{p-k}b^k + b^p.
		\]
		Проанализируем коэффициенты $\binom{p}{k} = \frac{p!}{k!(p-k)!}$. Так как $0 < k < p$, ни факториал $k!$, ни факториал $(p-k)!$ не содержат в своем разложении простого множителя $p$. Следовательно, число $p$ из числителя ни с чем не сокращается, и каждый такой биномиальный коэффициент можно представить в виде $p \cdot m$ для некоторого целого $m$.
		
		Так как в кольце $R$ любое кратное $p$ обращается в нуль ($p \cdot c = 0$), то и все промежуточные слагаемые суммы обращаются в нуль:
		\[
		\binom{p}{k} a^{p-k}b^k = p \cdot m \cdot (a^{p-k}b^k) = 0.
		\]
		В итоге остаются только крайние степени: $(a+b)^p = a^p + b^p$.
	\end{solution}
	
	\subsection*{Задача 4}
	Докажите, что многочлены
	\[
	4x^3 - 6x^2 + 3x - 6 \quad\text{и}\quad x^n - 2
	\]
	неприводимы над $\mathbb Q$ при любом натуральном $n$.
	
	\begin{solution}
		Для строгого доказательства неприводимости многочленов с целыми коэффициентами над полем рациональных чисел воспользуемся \textbf{критерием Эйзенштейна}. Критерий гласит: если для многочлена существует простое число $p$ такое, что старший коэффициент не делится на $p$, все остальные коэффициенты делятся на $p$, а свободный член делится на $p$, но не делится на $p^2$, то многочлен неприводим над $\mathbb{Q}$.
		
		\textbf{Первый многочлен:} $f(x)=4x^3 - 6x^2 + 3x - 6$. Применим критерий для простого числа $p=3$:
		\begin{itemize}
			\item старший коэффициент $a_3 = 4$ \textit{не делится} на $3$;
			\item промежуточные коэффициенты $a_2 = -6$ и $a_1 = 3$ \textit{делятся} на $3$;
			\item свободный член $a_0 = -6$ \textit{делится} на $3$, но \textit{не делится} на $p^2 = 9$.
		\end{itemize}
		Все три условия выполнены, следовательно, $f(x)$ неприводим над $\mathbb Q$.
		
		\textbf{Второй многочлен:} $g(x)=x^n - 2$. Применим критерий для простого числа $p=2$:
		\begin{itemize}
			\item старший коэффициент $a_n = 1$ \textit{не делится} на $2$;
			\item все промежуточные коэффициенты $a_{n-1} = \dots = a_1 = 0$ тривиально \textit{делятся} на $2$;
			\item свободный член $a_0 = -2$ \textit{делится} на $2$, но \textit{не делится} на $p^2 = 4$.
		\end{itemize}
		Критерий Эйзенштейна выполнен при абсолютно любом натуральном $n\ge 1$, поэтому многочлен $x^n - 2$ гарантированно неприводим над $\mathbb Q$.
	\end{solution}
	
	\subsection*{Задача 5}
	Вычислите дискриминант многочлена
	\[
	n!\left(1 + \frac{x}{1!} + \frac{x^2}{2!} + \dots + \frac{x^n}{n!}\right),
	\]
	используя связь дискриминанта и результанта.
	
	\begin{solution}
		Обозначим исходный многочлен как $f(x)$:
		\[
		f(x) = n! + n!x + \frac{n!}{2}x^2 + \dots + x^n.
		\]
		Это унитарный многочлен степени $n$ (старший коэффициент $a_n = 1$). Дискриминант унитарного многочлена выражается через результант самого многочлена и его производной строгой формулой:
		\[
		D(f) = (-1)^{\frac{n(n-1)}2} \operatorname{Res}(f, f').
		\]
		Найдём производную многочлена $f(x)$:
		\[
		f'(x) = n!\left(1 + x + \frac{x^2}{2!} + \dots + \frac{x^{n-1}}{(n-1)!}\right).
		\]
		Легко заметить связь между функцией и её производной:
		\[
		f(x) = f'(x) + x^n.
		\]
		Воспользуемся фундаментальным свойством результанта $\operatorname{Res}(A + BC, C) = \operatorname{Res}(A, C)$, применяя его к паре $(f, f')$:
		\[
		\operatorname{Res}(f, f') = \operatorname{Res}(f' + x^n, f') = \operatorname{Res}(x^n, f').
		\]
		По определению результанта через корни, $\operatorname{Res}(P, Q) = a_m^k \prod_{i=1}^m Q(\alpha_i)$, где $\alpha_i$ — корни $P(x)$, а $a_m$ — старший коэффициент. В нашем случае многочлен $P(x) = x^n$ имеет старший коэффициент $1$ и единственный корень $0$ кратности $n$. Следовательно:
		\[
		\operatorname{Res}(x^n, f') = 1^{\deg f'} \prod_{i=1}^n f'(0) = (f'(0))^n.
		\]
		Значение производной в нуле равно её свободному члену: $f'(0) = n! \cdot 1 = n!$. Подставляем это в результант:
		\[
		\operatorname{Res}(f, f') = (n!)^n.
		\]
		Итоговая формула дискриминанта принимает вид:
		\[
		D(f) = (-1)^{\frac{n(n-1)}2}\, (n!)^n.
		\]
	\end{solution}
	
\end{document}